\chapter{Time Dilation}

\section*{Introduction}

GPS satellites must correct for time dilation. Their clocks run about 38 microseconds per day slower than ground clocks due to their high speed (special relativity) and faster due to weaker gravity (general relativity). Without correction, GPS positions would drift by about 10 km per day. Time dilation is not science fiction---it is engineering reality.

\begin{enumerate}
\item \textbf{Special Relativistic Time Dilation:} A moving clock runs slow by the Lorentz factor.

\item \textbf{Gravitational Time Dilation:} Clocks in stronger gravitational fields run slower.

\end{enumerate}


\[
\Delta t = \gamma \Delta t_0 = \frac{\Delta t_0}{\sqrt{1 - v^2/c^2}}
\]

\section*{Fundamental Equations Used}

The derivation starts from the light-clock thought experiment: a photon bouncing between two parallel mirrors in a moving reference frame.

\section*{Assumptions}

\textbf{Assumptions:} The two clocks are in uniform relative motion (inertial frames). The speed of light is the same in all frames. No gravitational field (or identical gravitational potentials).


\textbf{Limitations:} Does not apply to accelerating clocks (need general relativity). The effect is negligible for $v \ll c$. Requires synchronization of clocks, which itself depends on the relativity of simultaneity.


\section*{Mathematical Derivation}

\textbf{Step 1: The Light Clock Thought Experiment.}

Imagine a ``light clock'': a photon bounces vertically between two mirrors separated by distance $L$. In the rest frame of the clock, one tick corresponds to the photon traveling from the bottom mirror to the top and back, a round-trip distance of $2L$. The time for one tick is:

\begin{equation}
\Delta t_0 = \frac{2L}{c}.
\end{equation}

This is the \textbf{proper time} $\Delta t_0$---the time measured by a clock at rest relative to the events it measures.

\textbf{Step 2: The Moving Light Clock.}

Now observe the same clock from a frame in which it moves horizontally at velocity $v$. During one tick, the clock has moved a horizontal distance $v\,\Delta t$. The photon now travels along a diagonal path (a ``zigzag''). The round-trip path consists of two diagonal legs, each of length $\sqrt{L^2 + (v\,\Delta t/2)^2}$, giving a total path length:

\begin{equation}
d = 2\sqrt{L^2 + \left(\frac{v\,\Delta t}{2}\right)^2}.
\end{equation}

\textbf{Step 3: Applying the Constancy of the Speed of Light.}

By the second postulate of special relativity, the speed of light is $c$ in all inertial frames. Therefore:

\begin{equation}
d = c\,\Delta t.
\end{equation}

Substituting:

\begin{equation}
c\,\Delta t = 2\sqrt{L^2 + \frac{v^2\,(\Delta t)^2}{4}}.
\end{equation}

\textbf{Step 4: Solving for $\Delta t$.}

Square both sides:

\begin{equation}
c^2\,(\Delta t)^2 = 4L^2 + v^2\,(\Delta t)^2.
\end{equation}

Collect terms involving $(\Delta t)^2$:

\begin{equation}
(c^2 - v^2)\,(\Delta t)^2 = 4L^2.
\end{equation}

Solve for $\Delta t$:

\begin{equation}
(\Delta t)^2 = \frac{4L^2}{c^2 - v^2} = \frac{4L^2}{c^2\left(1 - v^2/c^2\right)}.
\end{equation}

\begin{equation}
\Delta t = \frac{2L}{c\,\sqrt{1 - v^2/c^2}}.
\end{equation}

\textbf{Step 5: Expressing in Terms of Proper Time.}

Recall that $\Delta t_0 = 2L/c$. Substituting:

\begin{equation}
\Delta t = \frac{\Delta t_0}{\sqrt{1 - v^2/c^2}}.
\end{equation}

Defining the Lorentz factor $\gamma = 1/\sqrt{1 - v^2/c^2}$:

\begin{equation}
\boxed{\Delta t = \gamma\,\Delta t_0 = \frac{\Delta t_0}{\sqrt{1 - v^2/c^2}}.}
\end{equation}

Since $\gamma \geq 1$ for all $v < c$, the moving clock always runs slower: $\Delta t \geq \Delta t_0$. This is \textbf{special relativistic time dilation}.

\textbf{Step 6: Connection to the Lorentz Transformation.}

Time dilation emerges directly from the Lorentz transformation. If two events occur at the same location in frame $S'$ (a clock at rest), they are separated by $\Delta t_0$ in $S'$. In frame $S$, where $S'$ moves at velocity $v$:

\begin{equation}
\Delta t = \gamma\left(\Delta t' + \frac{v\,\Delta x'}{c^2}\right) = \gamma\,\Delta t',
\end{equation}

since $\Delta x' = 0$ (same location in $S'$). This confirms the result from the light clock and shows that time dilation is a direct consequence of the Lorentz structure of spacetime.

\section*{Characteristics of the Equation}

For $v \ll c$: $\gamma \approx 1 + v^2/2c^2$. At 10\% of $c$, time slows by 0.5\%. At 99\% of $c$, time slows by a factor of 7. The effect is reciprocal---each observer sees the other's clock as slow (resolved by the twin paradox).
\section*{Solution Methods}

Direct application of the formula for uniform motion. For accelerating frames: integrate proper time along the worldline. For gravitational fields: use the metric tensor and proper time integral.
\section*{Verifications}

Hafele-Keating experiment (1971): atomic clocks flown around the world matched predictions. Muon lifetime: atmospheric muons live longer than lab muons due to their high speed. Pound-Rebka experiment (1959): measured gravitational frequency shift in a tower.
\section*{Applications}

GPS satellite corrections, particle physics (muon lifetime extension), nuclear physics (lifetime of unstable particles in accelerators), astrophysics (gravitational redshift measurements), science fiction (interstellar travel).
\section*{What Richard Feynman Would Have Asked}

\subsection*{1. The Plain-English Translation}
\textit{``What does it really mean for time to run slower?''}

The Core Meaning: It means that the ``ticking rate'' of any physical process---atomic vibrations, biological aging, radioactive decay---slows down when that process is moving relative to an observer. This is not a mechanical effect on clocks; it is a fundamental property of spacetime itself. Moving clocks don't just measure time differently---they \emph{experience} time differently.

The Metaphor: Imagine two identical metronomes. Place one on a train moving at constant speed and one on the platform. From the platform, the moving metronome ticks more slowly. But from inside the train, the moving metronome ticks normally---it's the platform metronome that appears slow. Neither is ``wrong''; they are in different reference frames.

\subsection*{2. The Localized Mechanics}
\textit{``At the atomic level, what makes the clock run slower?''}

He would press you on whether time dilation is ``real'' or ``apparent.'' At the atomic level, a clock's ticking rate is determined by the frequency of electromagnetic radiation emitted by its atoms. Time dilation means that this radiation has a lower frequency (longer wavelength) when the clock is moving---the light itself is redshifted. This is not a mechanical stretching of the clock; it is a stretching of time itself.

The Smoothness Check: He would ask whether the astronaut feels anything different. No---all biological processes slow down equally. The astronaut's heartbeat, brain waves, and thought processes all slow by the same factor. Subjectively, time feels normal. The difference only becomes apparent upon comparison with a stationary clock.

\subsection*{3. Testing the Extreme Limits}
\textit{``What happens at speeds close to $c$?''}

The Speed-of-Light Limit: As $v \to c$, $\gamma \to \infty$. Time effectively stops for the moving object as seen by a stationary observer. No massive object can reach $c$---it would require infinite energy. This is why $c$ is the ultimate speed limit.

The Cosmic Ray Limit: Cosmic ray muons travel at 0.998$c$ and have a lab-frame lifetime 15.7 times longer than their rest-frame lifetime. Without time dilation, almost none would reach the Earth's surface---but millions do, exactly as predicted.

The Twin Paradox Resolution: One twin stays home, the other travels and returns. The traveling twin ages less. The asymmetry is that the traveling twin must accelerate (turn around), breaking the symmetry. The resolution is in general relativity: acceleration is equivalent to a gravitational field, which causes additional time dilation.

\subsection*{4. Dimensionality and Geometry}
\textit{``How does the geometry of spacetime explain time dilation?''}

He would point out that time dilation is a geometric effect in Minkowski spacetime. The ``distance'' in spacetime (the spacetime interval) is invariant, but the time and space components are not. Different observers ``slice'' spacetime differently---what one observer calls ``time,'' another calls a mixture of ``time'' and ``space.'' The Lorentz transformation is a rotation in spacetime.

\subsection*{5. Cross-Disciplinary Connection}
\textit{``Where else does this same relativistic effect appear?''}

Time dilation appears in: particle physics (particles in accelerators live longer), nuclear weapons (fission products moving at high speed decay slower), medical imaging (PET scanner timing relies on relativistic corrections), and even in financial modeling (``Black-Scholes'' equations have analogues in relativistic finance where volatility plays the role of velocity).
\section*{FAQ}

\textbf{Q: Does time dilation mean the moving person ages less?}\\
\textbf{A:} Yes. An astronaut traveling at 99\% of $c$ for 10 years (ship time) would return to find about 71 years have passed on Earth. This is not a paradox---it follows directly from the constancy of the speed of light.

\textbf{Q: Why doesn't the astronaut feel time passing slowly?}\\
\textbf{A:} Time dilation is only apparent when comparing clocks in different frames. In your own frame, time always passes normally. You only notice the difference when you compare with someone in a different frame.
\section*{Important Bibliography}

\begin{enumerate}
\item Einstein, A., ``Zur Elektrodynamik bewegter K\"orper,'' \textit{Annalen der Physik} \textbf{17}, 891--921 (1905).
\item Taylor, E.F. and Wheeler, J.A., \textit{Spacetime Physics}, 2nd ed. (W.H. Freeman, 1992), Chapters 1--2.
\item Hafele, J.C. and Keating, R.E., ``Around-the-World Atomic Clocks: Predicted Relativistic Time Gains,'' \textit{Science} \textbf{177}, 166--168 (1972).
\item Resnick, R., \textit{Introduction to Special Relativity} (Wiley, 1968), Chapters 1--4.
\item Mould, R.A., \textit{Special Relativity} (Springer, 1994), Chapter 3.
\end{enumerate}

